\clearpage
\setcounter{page}{1}
\maketitlesupplementary


\noindent\textbf{Semantic-Adaptive Segmentation.} Given a long video $\mathcal{X}$, we uniformly sample it with a fixed stride $\Delta t$ to obtain a sparse frame set $\mathcal{F} = \{ f_t \}_{t=1}^{N} \subset \mathcal{X}$. Each frame $f_i$ is processed by a language-aligned visual encoder $E_{\mathrm{vis}}$, such as the vision tower of an MLLM, to extract semantic embeddings $ \mathcal{Z}_{\text{feat}} = \{ z_t \}_{t=1}^{N} = E_{\mathrm{vis}}(\mathcal{F})$ of the sampled frames.

To adaptively segment the video based on semantic continuity, we reformulate the problem as an optimal partitioning task, solved via semantic change point detection along the temporal axis. Given the frame-level semantic embeddings $\mathcal{Z}_\text{feat}$, our goal is to identify an optimal set of temporal boundaries $\tau = \{\tau_k^{\text{start}}, \tau_k^{\text{end}}\}_{k=1}^{K}$ that segments the video into $K$ semantically homogeneous short clips. To pinpoint these boundaries effectively, we quantify the magnitude of frame-to-frame semantic divergence to detect discontinuities in the feature stream. Concretely, we first derive the semantic flux sequence $\mathcal{D}=\{d_t\}_{t=2}^N$, where $d_t = 1 - \frac{z_t \cdot z_{t-1}}{\|z_t\|\|z_{t-1}\|}$ captures the instantaneous semantic shift. We then employ the Pruned Exact Linear Time (PELT) algorithm to solve the following optimal partitioning objective:
\begin{equation}
\min_{\{ \tau_k^{\text{start}}, \tau_k^{\text{end}} \}}
\sum_{k=1}^K \mathcal{C}_{\text{ker}}(d_{\tau_k^{\text{start}} : \tau_k^{\text{end}}})
+ K \cdot \beta,
\end{equation}
where a kernel-based cost function $\mathcal{C}_{\text{ker}}$ quantifies the semantic inconsistency of the segment over the semantic flux, and the penalty $K\cdot\beta$ prevents over-segmentation by discouraging unnecessary semantic change points. We provide a more detailed explanation of the cost function $\mathcal{C}_{\text{ker}}$ in Appendix~\ref{appendix:cost_function}. Once the optimal temporal boundaries $\tau$ and the number of clips $K$ are determined, we extract segmented clips from the dense video stream as $S = \{\, \{\, x_t \mid t \in [\tau_k^{\text{start}}, \tau_k^{\text{end}}] \,\} \subset \mathcal{X} \,\}_{k=1}^{K}$. This process adaptively segments the long video into short clips that are both temporally continuous and semantically aligned, serving as the node candidates for constructing the subsequent spatio-temporal graph.
